\section{Fungsi: Deklarasi, Parameter, dan Return}

Fungsi adalah blok kode yang dapat dipanggil berulang kali dengan nama tertentu. Fungsi menerima \textbf{parameter} (input), melakukan operasi, dan dapat mengembalikan nilai dengan \code{return}. Fungsi adalah fondasi organisasi kode: memecah program besar menjadi bagian-bagian kecil yang dapat diuji dan digunakan ulang. Di JavaScript, fungsi adalah \textit{first-class citizens}---dapat disimpan di variabel, dikirim sebagai argumen, dan dikembalikan dari fungsi lain \cite{jsInfo}.

\begin{webcode}[language=JavaScript, caption={Tiga cara membuat fungsi}]
// 1. Function declaration (hoisted)
function greet(name) {
  return `Hello, ${name}!`;
}

// 2. Function expression (tidak hoisted)
const add = function(a, b) {
  return a + b;
};

// 3. Arrow function (ES6, ringkas)
const multiply = (a, b) => a * b;

// Default parameter
function power(base, exp = 2) {
  return base ** exp;
}
console.log(power(3));    // 9 (exp default 2)
console.log(power(3, 3)); // 27
\end{webcode}

\begin{table}[htbp]
\centering
\begin{tabular}{|P{3.5cm}|P{4cm}|P{4.5cm}|}
\hline
\textbf{Jenis} & \textbf{Sintaks} & \textbf{Karakteristik} \\
\hline
Declaration & \code{function f() \{\}} & Hoisted; dapat dipanggil sebelum dideklarasikan \\
\hline
Expression & \code{const f = function() \{\}} & Tidak hoisted; di variabel \\
\hline
Arrow (ES6) & \code{const f = () => \{\}} & Ringkas; tanpa \code{this} sendiri \\
\hline
\end{tabular}
\caption{Tiga cara mendefinisikan fungsi JavaScript}
\end{table}

\begin{konsep}
\textbf{Arrow Function dan \code{this}.} Arrow function tidak memiliki \code{this} sendiri---ia mewarisi \code{this} dari konteks luar (\textit{lexical this}). Ini sangat berguna dalam callback dan event handler di dalam objek. Namun, arrow function \textbf{tidak cocok} untuk method objek (\code{this} tidak merujuk ke objek) dan tidak bisa digunakan sebagai constructor. Gunakan arrow function untuk callback dan function expression biasa untuk method objek \cite{mdnJsGuide}.
\end{konsep}

